\chapter{Fuentes de datos, integración e ingeniería de características}\label{cap:datos_features}

Este capítulo documenta la cadena que transforma tres ficheros crudos, heterogéneos en
formato y en semántica, en la matriz de 161 variables sobre la que se entrenan los modelos
de los capítulos \ref{cap:metodologia} y \ref{cap:resultados}. Cada etapa hace explícita
la regla anti-fuga que la gobierna, puesto que la integridad de esa cadena es la condición
de validez de todo resultado reportado.

\section{Fuentes}\label{sec:fuentes_datos}

El proyecto integra tres fuentes crudas, todas de frecuencia diaria y con un horizonte
temporal común de 1.310 observaciones continuas, sin huecos y sin valores nulos, entre el
1 de enero de 2023 y el 2 de agosto de 2026. Los tres ficheros se conservan sin modificar
en \texttt{data/\allowbreak{}raw/} y se resumen en la tabla \ref{tab:fuentes_datos}.

La primera fuente procede del Consorcio Regional de Transportes de Madrid (CRTM), el
organismo público responsable de la coordinación tarifaria e informativa del transporte
colectivo en la Comunidad de Madrid, y reporta la evolución diaria de validaciones
agregadas por operador: Metro de Madrid, EMT (autobús urbano), las concesionarias
interurbanas por carretera y Renfe Cercanías, más el total sumado de los cuatro. Se
distribuye como un libro Excel de dos hojas (\texttt{mensual} y \texttt{diaria}, de la que
solo se emplea la segunda) cuyo maquetado difiere del descrito en la documentación inicial
del proyecto: la cabecera no ocupa la primera fila, sino la segunda, de modo que una
lectura ingenua con \texttt{header=2} en lugar de \texttt{header=1} consume
silenciosamente la primera observación (2023-01-01) como si fuera cabecera.

La columna de fechas, además, carece de nombre propio en el fichero de origen, la primera
columna contiene un marcador suelto sin significado analítico, la octava está vacía y la
novena recoge una nota que advierte del carácter \emph{provisional} de las cifras de
demanda; esa advertencia se traslada íntegra a las limitaciones del proyecto
(\ref{sec:limitaciones}). Las cuatro columnas de operador se conservan explícitamente en
el conjunto unificado y en la matriz de características, aunque el objetivo de modelado
sea exclusivamente \texttt{total}: no son variables redundantes a descartar, sino una
desagregación que sostiene el análisis exploratorio de la sección \ref{sec:eda} y podría
alimentar líneas futuras sobre demanda por modo de transporte.

La segunda fuente es una serie meteorológica diaria de Open-Meteo para la localización de
Madrid (40,386642° N, 3,676086° O, 666 m de altitud), con 22 variables que cubren
temperatura, temperatura aparente, precipitación, humedad relativa, viento, presión a
nivel del mar, radiación de onda corta y código de tiempo (WMO). La propuesta inicial
preveía recurrir a la Agencia Estatal de Meteorología (AEMET); se optó finalmente por
Open-Meteo porque su producto de reanálisis cubre la ventana completa 2023-2026 para un
único punto geográfico sin huecos ni cambios de estación de medida, lo que evita imputar o
empalmar series. Esa operación, sobre variables que después se retardan y promedian,
habría introducido una fuente de error difícil de acotar. La misma familia de API se
retoma en la sección \ref{sec:trabajos_futuros} como vía natural para incorporar
pronóstico operativo.

El fichero incorpora dos filas de metadatos de cabecera y una línea en blanco antes de la
fila con los nombres de columna reales, y esos metadatos declaran una zona horaria
\texttt{Europe/\allowbreak{}Berlin} en lugar de \texttt{Europe/\allowbreak{}Madrid}; ambas comparten horario de
verano centroeuropeo (CEST) en la práctica totalidad del año, por lo que el efecto sobre
los agregados diarios se consideró y descartó como fuente de sesgo relevante durante la
ingesta, sin darlo por supuesto a priori. Las unidades viajan incrustadas en el nombre de
cada columna (por ejemplo, \texttt{temperature\_2m\_mean (°C)} o
\texttt{shortwave\_radiation\_sum (MJ/\allowbreak{}m²)}) y se normalizan mediante una expresión
regular que las separa del nombre canónico.

\begin{sloppypar}
La tercera fuente es el calendario laboral de referencia para la Comunidad de Madrid, un
CSV separado por punto y coma y codificado en UTF-8 con marca de orden de bytes (BOM), con
las fechas en formato \texttt{dd/\allowbreak{}mm/\allowbreak{}aaaa} y analizadas mediante un patrón explícito en
lugar de una inferencia automática de orden día/mes, para evitar ambigüedades
sistemáticas. Distingue el día de la semana (en minúsculas y sin tildes, de modo que
\texttt{miercoles} o \texttt{sabado} no deben compararse contra sus formas acentuadas), el
tipo de día (\texttt{laborable}/\texttt{sabado}/\texttt{domingo}/\texttt{festivo}) y,
cuando procede, el tipo y el nombre de la festividad. El nombre de la columna de tipo de
día declara una quinta categoría, \texttt{domingo festivo}, que no se observó en ninguna
de las 1.310 filas disponibles; esta ausencia se verificó de forma explícita durante la
ingesta, en lugar de asumirse por la mera lectura de esa documentación, y la categoría se
mantiene declarada en la codificación para que, si apareciera en una actualización futura
del calendario, se codifique correctamente en vez de degradarse silenciosamente a un valor
nulo.
\end{sloppypar}

\begin{table}[htbp]
  \centering
  \small
  \setlength{\tabcolsep}{4pt}
  \resizebox{\textwidth}{!}{\begin{tabular}{>{\raggedright\arraybackslash}p{3.5cm}>{\raggedright\arraybackslash}p{4.6cm}>{\raggedright\arraybackslash}p{2.8cm}>{\raggedright\arraybackslash}p{3.6cm}}
\toprule
\rowcolor{naranja}
\textbf{Fichero} & \textbf{Contenido} & \textbf{Formato} & \textbf{Origen} \\
\midrule
\rowcolor{filaclara} CRTM\_\allowbreak{}Evolucion\_\allowbreak{}demanda\_\allowbreak{}diaria.xlsx & Demanda diaria por operador (metro, EMT, carretera, cercanías) y total & Excel, hoja 'diaria', cabecera en fila 2 & \url{https://datos.crtm.es/documents/crtm::crtm-evolucion-demanda-diaria/about} \\
open-meteo-40.39N3.68W666m.csv & 22 variables meteorológicas diarias (temperatura, precipitación, viento, presión, radiación) & CSV, cabecera real en línea 4 & \url{https://open-meteo.com/en/docs/historical-weather-api} \\
\rowcolor{filaclara} 300082-1-calendario\_\allowbreak{}laboral-csv.csv & Día de la semana, tipo de día (laborable/\allowbreak{}sábado/\allowbreak{}domingo/\allowbreak{}festivo) y festividad & CSV separado por ';', UTF-8 con BOM & \url{https://datos.madrid.es/dataset/300082-0-calendario_laboral} \\
\bottomrule
\end{tabular}
}
  \caption{\headlinecolor{Fuentes de datos crudas (\texttt{data/\allowbreak{}raw/}), 2023-01-01 a 2026-08-02.}}
  \label{tab:fuentes_datos}
\end{table}

La figura \ref{fig:serie_completa_total} muestra la serie completa de demanda total tal
como resulta de esta primera fuente, antes de cualquier transformación: fija la escala del
problema (entre 1,3 y 6,6 millones de viajes diarios aproximadamente) y deja ya visible la
textura semanal de sierra y la caída estacional recurrente cada agosto, patrones que la
sección \ref{sec:eda} analiza en detalle.

\begin{figure}[htbp]
  \centering
  \includegraphics[width=0.78\textwidth]{fig_serie_completa_total.png}
  \caption{\headlinecolor{Demanda diaria total de transporte público en Madrid, serie completa.}}
  \label{fig:serie_completa_total}
\end{figure}

\section{Unificación y validación}\label{sec:unificacion_validacion}

Las tres fuentes se combinan mediante una unión interna (\emph{inner join}) sobre la
columna de fecha. Se prefiere a una unión externa por la izquierda porque esta fabricaría
filas exógenas enteramente nulas allí donde una fuente tuviera un hueco, contaminando
artificialmente el indicador de calidad ``cero nulos'' del conjunto unificado. La
contrapartida es que una unión interna oculta, por construcción, cualquier hueco de
cobertura: si una fecha faltara simultáneamente en las tres fuentes, el problema no se
manifestaría como una fila nula sino como una fila que no llega a existir. Por ello la
unificación aplica cuatro comprobaciones de integridad temporal en cadena: fechas
duplicadas dentro de cada fuente, antes de fusionar nada; el número de filas resultante
contra el de \emph{cada} fuente por separado, interrumpiendo la ejecución ante cualquier
discrepancia y señalando explícitamente las fechas concretas que se perderían, en lugar de
un recuento agregado; una reindexación del resultado contra un rango de fechas diario
completo entre el mínimo y el máximo observados, única comprobación capaz de detectar un
hueco común a las tres fuentes simultáneamente, precisamente el caso que la comparación
anterior no puede ver; y una verificación final de ausencia de duplicados y de valores
nulos sobre el conjunto ya unificado.

Las cuatro se someten además a pruebas negativas que inyectan huecos sintéticos en
conjuntos de datos de prueba en memoria y confirman que cada una interrumpe la ejecución
cuando corresponde, de manera que no puedan degradarse silenciosamente en verificaciones
inertes con el tiempo.

El resultado se persiste en \texttt{data/\allowbreak{}interim/\allowbreak{}unified\_daily.parquet} como una tabla
de 1.310 filas por 31 columnas, resumidas por grupos temáticos en la tabla
\ref{tab:esquema_unificado}: la fecha, las cinco columnas de demanda (el objetivo
\texttt{total} y los cuatro operadores), las veintiuna variables meteorológicas y las
cuatro columnas de calendario. Sobre ella se verifica, y se mantiene como invariante
durante todo el proyecto, que la suma de las cuatro columnas de operador reproduce
exactamente \texttt{total} en las 1.310 filas, salvo por una tolerancia de $10^{-6}$ que
absorbe el ruido de redondeo de las fórmulas de Excel del fichero de origen (por ejemplo,
un valor de \texttt{carretera} igual a $847904{,}00000001$ en la fila del 7 de julio de
2026, con una desviación del orden de $10^{-8}$ respecto al entero esperado); cualquier
desviación genuinamente fraccionaria, por encima de esa tolerancia, se trata como un error
de datos y detiene el proceso, ya que ensancharla para acallar un futuro fallo eliminaría
la capacidad de detectar una ruptura real de esa identidad.

Por último, los nulos que el calendario de origen deja en \texttt{holiday\_type} y
\texttt{holiday\_name} para los días sin festividad se sustituyen deliberadamente por el
centinela \texttt{'none'}: así la invariante de ``cero nulos'' conserva su significado
—cualquier nulo que aparezca en adelante señala un defecto real, no la ausencia ordinaria
de festividad— y la ingeniería de características recibe un nivel categórico limpio en
lugar de un hueco que codificar de forma ad hoc.

\begin{table}[htbp]
  \centering
  \small
  \begin{tabular}{>{\raggedright\arraybackslash}p{3.0cm}>{\centering\arraybackslash}p{2.0cm}>{\raggedright\arraybackslash}p{9.0cm}}
\toprule
\rowcolor{naranja}
\textbf{Grupo} & \textbf{N.º columnas} & \textbf{Ejemplo} \\
\midrule
\rowcolor{filaclara} Fecha & 1 & date \\
Demanda (target + operadores) & 5 & total, metro, emt, carretera, cercanias \\
\rowcolor{filaclara} Meteorología & 21 & temperature\_\allowbreak{}2m\_\allowbreak{}mean, precipitation\_\allowbreak{}sum, wind\_\allowbreak{}speed\_\allowbreak{}10m\_\allowbreak{}max \\
Calendario & 4 & day\_\allowbreak{}of\_\allowbreak{}week\_\allowbreak{}es, day\_\allowbreak{}type, holiday\_\allowbreak{}type, holiday\_\allowbreak{}name \\
\bottomrule
\end{tabular}

  \caption{\headlinecolor{Esquema unificado (\texttt{data/\allowbreak{}interim/\allowbreak{}unified\_daily.parquet}), 1.310 x 31.}}
  \label{tab:esquema_unificado}
\end{table}

La figura \ref{fig:demanda_por_operador} descompone la serie total por operador y hace
visible su jerarquía de escala: Metro de Madrid concentra el volumen mayor, seguido de
EMT, mientras que las concesionarias por carretera y Renfe Cercanías aportan magnitudes
sensiblemente menores pero con una textura semanal igualmente marcada. No participa en el
modelado del capítulo \ref{cap:metodologia}, cuyo único objetivo es \texttt{total}, pero
justifica que las cuatro columnas de operador se conserven íntegras a través de todo el
pipeline: son la base de este análisis exploratorio secundario y quedan disponibles para
un trabajo futuro de predicción desagregada por modo.

\begin{figure}[htbp]
  \centering
  \includegraphics[width=0.78\textwidth]{fig_demanda_por_operador.png}
  \caption{\headlinecolor{Demanda diaria por operador del CRTM (Metro, EMT, carretera, Cercanías).}}
  \label{fig:demanda_por_operador}
\end{figure}

\section{Análisis exploratorio de datos}\label{sec:eda}

El análisis exploratorio persigue un objetivo acotado: caracterizar la estructura de
estacionalidad y de dispersión de la demanda antes de diseñar sobre ella la ingeniería de
características de la sección \ref{sec:ingenieria_features}, no agotar todas las lecturas
posibles de los datos. La serie completa (figura \ref{fig:serie_completa_total}) ya deja
ver, a simple vista y de forma repetida en cada uno de los tres años y medio cubiertos,
dos regularidades anuales: un desplome pronunciado y sostenido en agosto —el periodo
vacacional de mayor intensidad en Madrid, en el que se suspende buena parte de la
actividad laboral y educativa que sostiene la demanda de días laborables— y una
recuperación progresiva en septiembre que culmina en un máximo relativo hacia octubre y
noviembre, una vez restablecida por completo la actividad laboral y escolar.

La figura \ref{fig:estacionalidad_anual} lo cuantifica: la demanda media de agosto ronda
los 3 millones de viajes diarios, sensiblemente por debajo de cualquier otro mes, mientras
que octubre alcanza el máximo anual con una media próxima a los 5,2 millones; invierno y
primavera se mantienen en una meseta intermedia, entre 4,5 y 5 millones, con una
dispersión (barras de error de desviación típica) sistemáticamente mayor en los meses que
combinan más irregularidades de calendario —festivos nacionales y autonómicos, semana de
Navidad, Semana Santa de fecha variable— que en los meses centrales del curso.

\begin{figure}[htbp]
  \centering
  \includegraphics[width=0.78\textwidth]{fig_estacionalidad_semanal.png}
  \caption{\headlinecolor{Estacionalidad semanal de la demanda total (diagrama de caja por día de la semana).}}
  \label{fig:estacionalidad_semanal}
\end{figure}

La segunda regularidad, de periodicidad semanal y de magnitud mayor que la anual, está en
la figura \ref{fig:estacionalidad_semanal}. Los cinco días laborables presentan medianas
muy próximas entre sí, entre 5 y 5,5 millones de viajes diarios, con una cola inferior de
valores atípicos que corresponde a los festivos caídos en día laborable; el sábado
desciende a una mediana cercana a los 3,4 millones y el domingo, el día de menor demanda,
se sitúa alrededor de los 2,6 millones. La transición entre régimen laborable y fin de
semana implica, por tanto, un salto de aproximadamente 2,5 millones de viajes diarios dos
veces por semana, magnitud clave para interpretar correctamente los baselines del capítulo
\ref{cap:metodologia}: un modelo de persistencia ingenuo ($\hat y_t = y_{t-1}$) queda
sistemáticamente mal calibrado en cada transición viernes-sábado y domingo-lunes,
precisamente porque ignora esta estructura semanal tan pronunciada.

\begin{figure}[htbp]
  \centering
  \includegraphics[width=0.78\textwidth]{fig_estacionalidad_anual.png}
  \caption{\headlinecolor{Estacionalidad anual de la demanda total (media mensual, 2023-2026).}}
  \label{fig:estacionalidad_anual}
\end{figure}

Por último, la figura \ref{fig:meteo_vs_demanda} cruza la demanda total con la
temperatura media diaria, coloreando cada observación por tipo de día. Es informativa
precisamente por lo que no muestra: dentro de cada franja de tipo de día la nube de puntos
se mantiene aproximadamente horizontal, sin una relación monótona clara entre temperatura
y demanda, mientras que la separación vertical entre franjas (laborable frente a sábado,
domingo y festivo) domina visualmente el gráfico. Sugiere que el \textbf{calendario, no la
meteorología, es el factor exógeno de primer orden} sobre la demanda agregada, y que la
meteorología, si aporta señal, lo hace como efecto secundario que matiza el nivel de cada
régimen de calendario más que como determinante independiente. Esta lectura motiva la
decisión, desarrollada en la sección \ref{sec:promocion_meteo}, de incorporar la
meteorología como bloque de variables exógenas retardadas y nunca como sustituto del
calendario.

\begin{figure}[htbp]
  \centering
  \includegraphics[width=0.78\textwidth]{fig_meteo_vs_demanda.png}
  \caption{\headlinecolor{Demanda total frente a temperatura media diaria, por tipo de día (EDA; la meteorología del mismo día no entra en el modelo).}}
  \label{fig:meteo_vs_demanda}
\end{figure}

\section{Ingeniería de características}\label{sec:ingenieria_features}

A partir del esquema unificado de la sección \ref{sec:unificacion_validacion} se
construye una matriz de 161 variables de modelado, todas identificadas por el prefijo
\texttt{feat\_} y recuperables mecánicamente mediante \texttt{feature\_columns()} en lugar
de una lista mantenida a mano. Es una decisión deliberada para que el conjunto de
variables no pueda derivar silenciosamente respecto del código que las genera. La tabla
\ref{tab:grupos_features} resume los seis grupos resultantes, cuyo diseño matemático se
detalla a continuación.

El primer grupo son los retardos de la propia demanda total,
$\text{lag}_k(\mathit{total})_t = \mathit{total}_{t-k}$ para
$k \in \{1, 2, 3, 7, 14, 21, 28\}$: siete variables que cubren desde la inercia de un día
para otro hasta la memoria calendario de cuatro semanas, elegida porque el retardo 28
recoge el mismo día de la semana cuatro veces atrás, bajo el mismo régimen semanal que la
sección \ref{sec:eda} identificó como dominante. El mismo esquema se aplica a cada columna
de operador, 28 variables más: el reparto modal de ayer es información legítimamente
disponible en el momento de la predicción y conserva señal real, a diferencia del reparto
modal del propio día que se discute más abajo. Se suman cuatro estadísticos móviles de la
demanda total (media y desviación típica sobre ventanas de 7 y 28 días) y seis términos de
Fourier deterministas que codifican la estacionalidad semanal ($k=1$) y anual ($k=1$ y
$k=2$) mediante pares seno-coseno:
\[
  \sin\!\left(\frac{2\pi k t}{T}\right), \qquad
  \cos\!\left(\frac{2\pi k t}{T}\right),
\]
con $T=7$ para el término semanal y $T=365{,}25$ para los anuales, y $t$ medido en días
desde una época fija, \texttt{FOURIER\_EPOCH = 2023-01-01}, deliberadamente distinta de la
fecha mínima observada en cada ejecución: fijarla evita que la llegada de datos futuros
modifique retroactivamente la codificación de una fecha ya existente.

El último grupo interno, de once variables, cubre el calendario mediante codificación
binaria exhaustiva (\emph{one-hot encoding}) del tipo de día y del tipo de festividad,
incluyendo explícitamente el nivel centinela \texttt{'none'} y la categoría de cero
observaciones \texttt{domingo festivo}, de modo que una fila con todo ceros sería un error
de codificación detectable y no una ambigüedad entre ``categoría base'' y ``fallo
silencioso''. Se añaden un indicador binario de fin de semana y uno de ``día puente'': un
día laborable inmediatamente anterior o posterior a un festivo o domingo festivo,
construido observando también el calendario del día siguiente. Esta mirada hacia adelante
no constituye fuga de información porque el calendario oficial de festivos español se
conoce con más de un año de antelación y es, por tanto, dato disponible en el momento de
la predicción; el procedimiento detectó 55 días puente en todo el periodo, todos
laborables por definición.

\begin{table}[htbp]
  \centering
  \small
  \begin{tabular}{>{\raggedright\arraybackslash}p{4.3cm}>{\centering\arraybackslash}p{2.0cm}>{\raggedright\arraybackslash}p{7.8cm}}
\toprule
\rowcolor{naranja}
\textbf{Grupo} & \textbf{N.º variables} & \textbf{Descripción} \\
\midrule
\rowcolor{filaclara} Retardos de la demanda total & 7 & lags 1, 2, 3, 7, 14, 21, 28 \\
Retardos de operadores & 28 & 4 operadores x 7 retardos \\
\rowcolor{filaclara} Medias/\allowbreak{}desv. móviles de la demanda & 4 & media y desv. típica, ventanas 7 y 28, shift-first \\
Términos de Fourier & 6 & semanal k=1, anual k=1 y k=2 \\
\rowcolor{filaclara} Calendario & 11 & day\_\allowbreak{}type (5), holiday\_\allowbreak{}type (4), is\_\allowbreak{}weekend, is\_\allowbreak{}bridge\_\allowbreak{}day \\
Meteorología retardada & 105 & 21 variables x (lags 1/\allowbreak{}2/\allowbreak{}3/\allowbreak{}7 + media móvil 7) \\
\rowcolor{filaclara} Total & 161 &  \\
\bottomrule
\end{tabular}

  \caption{\headlinecolor{Grupos de variables del modelo, 161 en total (\texttt{feature\_columns()}).}}
  \label{tab:grupos_features}
\end{table}

Dos reglas anti-fuga gobiernan la construcción de este bloque y se verifican mediante
pruebas automatizadas específicas. La primera exige que todo retardo o estadístico móvil
derivado de \texttt{total} o de una columna de operador aplique un desplazamiento de un
día, \texttt{shift(1)}, \emph{antes} de cualquier operación de ventana móvil: es decir,
$\text{roll}_w(s.\text{shift}(1))_t$ y nunca $\text{roll}_w(s)_t.\text{shift}(1)$ como
corrección a posteriori. Para una ventana móvil retrospectiva (\emph{trailing}) ambos
órdenes son matemáticamente equivalentes, de modo que la regla no corrige aquí ningún
error numérico observado; se exige de todos modos porque su corrección no depende de esa
equivalencia particular, que deja de cumplirse en cuanto la ventana se centra
(\texttt{center=True}), y porque una ventana sin desplazar en absoluto sí constituye una
fuga real y directa: la media móvil en el instante $t$ incorporaría el propio valor $y_t$
que se pretende predecir.

La segunda es un guardián de fuga contemporánea sobre las columnas de operador: puesto que
$\mathit{metro}_t + \mathit{emt}_t + \mathit{carretera}_t + \mathit{cercanias}_t =
\mathit{total}_t$ de forma exacta (sección \ref{sec:unificacion_validacion}), cualquier
columna de operador del \emph{mismo} día no es una variable predictiva, sino una
reformulación aritmética del propio objetivo, y su presencia en la matriz haría que un
modelo alcanzara un ajuste casi perfecto en entrenamiento para colapsar en inferencia,
momento en el que el reparto modal del día en curso no puede conocerse todavía. Se aplica
dos veces: sobre la salida de cada módulo de retardos y de nuevo sobre el fichero final
ensamblado, que es el punto donde una columna mal etiquetada terminaría incorporándose a
la matriz.

Conviene decir con precisión qué se calcula antes de la partición y por qué es seguro.
Las 161 variables se construyen sobre las 1.310 filas completas, antes de que actúe
ninguna función de partición (\texttt{build\_features} precede a
\texttt{chronological\_split} en todos los flujos). No introduce fuga porque todas son
\emph{causales}: el valor de cualquier retardo, estadístico móvil desplazado o variable
meteorológica retardada en la fila $t$ depende únicamente de filas estrictamente
anteriores a $t$, y ninguna estima parámetro alguno a partir de los datos (una media móvil
desplazada es una función fija de sus $w$ valores previos, no un estadístico ajustado
sobre la serie). Los términos de Fourier y las variables de calendario son, además,
funciones cerradas de la fecha que no observan el objetivo en ningún momento. Calcular
cualquiera de estas variables sobre el rango completo es, por tanto, idéntico fila a fila
a calcularla por separado dentro de cada partición. La única transformación que sí se
\emph{ajusta} sobre datos es el escalador \texttt{MinMaxScaler} de la etapa 1, y se ajusta
exclusivamente sobre el bloque de entrenamiento, y dentro de cada fold del esquema OOF
sobre la porción de entrenamiento de ese fold (sección \ref{sec:esquema_oof}).

\subsection{Promoción meteorológica: solo variables retardadas}\label{sec:promocion_meteo}

La sección \ref{sec:eda} dejó abierta, más que resuelta, la pregunta de si la
meteorología debía incorporarse a la matriz observada en tiempo real o bajo alguna forma
de retardo. Este proyecto la resuelve de manera conservadora: las 21 variables
meteorológicas del esquema unificado se promocionan a variables de modelado exclusivamente
en forma retardada, combinando los retardos de 1, 2, 3 y 7 días con una media móvil de 7
días (también construida con desplazamiento previo, por la misma regla anti-fuga de la
sección anterior), lo que produce $21 \times 5 = 105$ variables nuevas y eleva la matriz
de 56 a 161 columnas. La meteorología del propio día de la predicción no entra nunca en el modelo, ni siquiera
junto a las retardadas, y la prohibición se verifica mediante un guardián dedicado
(\texttt{\_assert\_no\_same\_day\_weather}) que se reejecuta sobre el conjunto final
ensamblado. La razón es que la meteorología observada del mismo día no es, en rigor, un
dato disponible en el momento operativo de la predicción, sino una observación posterior a
esta; usarla sin retardo equivaldría a asumir de forma implícita un pronóstico
meteorológico perfecto, un sesgo optimista que además se concentraría precisamente en los
días atípicos (temporales, olas de calor) en los que un pronóstico real es menos fiable y
en los que el efecto sobre la demanda sería, presumiblemente, mayor.

El conjunto de retardos elegido para la meteorología, $\{1, 2, 3, 7\}$, es
deliberadamente más corto que el de la demanda total, $\{1, 2, 3, 7, 14, 21, 28\}$, porque
ambas series obedecen a mecanismos de memoria distintos: la autocorrelación de la demanda
a 28 días está gobernada por el calendario (el mismo día de la semana, cuatro semanas
atrás, bajo un régimen semanal recurrente), mientras que la autocorrelación atmosférica
carece de un mecanismo equivalente y decae en la práctica a lo largo de una escala
sinóptica de tres a siete días; extender los retardos meteorológicos a 14, 21 o 28 días
habría añadido 63 columnas prácticamente redundantes con los términos de Fourier anuales,
que ya capturan de forma más limpia la componente puramente estacional del clima. Queda como línea pendiente, y explícitamente fuera del alcance de los resultados
principales de este TFM (\ref{sec:limitaciones}), una variante de sensibilidad con
meteorología del mismo día sin retardar, que solo tendría sentido reportar como cota
superior bajo el supuesto idealizado de un pronóstico meteorológico perfecto, y nunca
como modelo de referencia.

\section{Particiones y fuga de datos}\label{sec:particiones_fuga}

Todos los modelos de los capítulos \ref{cap:metodologia} y \ref{cap:resultados}
(baselines estadísticos, LSTM, XGBoost y sus combinaciones) comparten una única partición
cronológica 70/15/15, calculada por una sola función, \texttt{chronological\_split}, que
actúa como fuente única de verdad: ningún modelo recalcula sus propias fronteras
temporales, precisamente para que un desacuerdo de un solo día entre dos modelos sobre
dónde empieza el test no pueda invalidar silenciosamente su comparación. La partición
resultante, resumida en la tabla \ref{tab:particiones}, reserva 917 días (70,0\,\%) para
entrenamiento, desde el 1 de enero de 2023 hasta el 5 de julio de 2025; 196 días
(15,0\,\%) para validación, hasta el 17 de enero de 2026; y 197 días (15,0\,\%) para test,
hasta el final de la serie disponible, el 2 de agosto de 2026.

El cálculo de estas fronteras exige dos precauciones que, de omitirse, introducirían
errores sutiles y difíciles de detectar a posteriori. La primera es numérica: $0{,}70
\times 1310$ evalúa en aritmética de punto flotante a $916{,}9999\ldots$ (esto es,
$916{,}\overline{9}$ hasta la precisión de la máquina), de modo que un truncamiento
ingenuo mediante \texttt{int()} daría 916 días de entrenamiento en lugar de los 917 que
resultan de $\lfloor 0{,}70 \times 1310 \rfloor$; la implementación añade una tolerancia
$\varepsilon = 10^{-9}$ antes de truncar, precisamente para blindarse frente a ese error
de redondeo binario. La segunda es de diseño: las fronteras se resuelven sobre fracciones
acumuladas ($\lfloor 0{,}70 \, n \rfloor$ para el final de entrenamiento, $\lfloor 0{,}85
\, n \rfloor$ para el de validación) en lugar de aplicar el porcentaje de forma
independiente a cada partición, de modo que los tres bloques encajen exactamente, sin
solapamientos ni huecos de un día que un redondeo partición por partición sí podría
introducir.

\begin{table}[htbp]
  \centering
  \small
  \begin{tabular}{lllrr}
\toprule
\rowcolor{naranja}
\textbf{Partición} & \textbf{Inicio} & \textbf{Fin} & \textbf{Días} & \textbf{Porcentaje} \\
\midrule
\rowcolor{filaclara} Train & 2023-01-01 & 2025-07-05 & 917 & 70,0 \\
Val & 2025-07-06 & 2026-01-17 & 196 & 15,0 \\
\rowcolor{filaclara} Test & 2026-01-18 & 2026-08-02 & 197 & 15,0 \\
\bottomrule
\end{tabular}

  \caption{\headlinecolor{Partición cronológica 70/15/15 (\texttt{src.utils.splits.chronological\_split}).}}
  \label{tab:particiones}
\end{table}

Sobre esta partición actúa, además, una política de calentamiento (\emph{warm-up})
derivada de la ingeniería de características de la sección \ref{sec:ingenieria_features}:
las 28 primeras filas de entrenamiento (2023-01-01 a 2023-01-28) contienen nulos en las
variables que dependen del retardo o la ventana móvil más larga, 28 días, y se descartan,
reduciéndolo de 917 a 889 filas efectivas; validación y test, muy por delante de esa
ventana inicial, no se ven afectados y conservan sus 196 y 197 días íntegros. La aplica
\texttt{apply\_warmup\_policy}, que \emph{interrumpe} la ejecución si el calentamiento
necesario llegara a alcanzar validación o test, en lugar de recortar silenciosamente el
periodo de evaluación: un conjunto de evaluación acortado sin aviso produciría métricas
que dejarían de ser comparables entre modelos y entre reejecuciones del pipeline, un
riesgo que se considera preferible convertir en un fallo explícito y temprano.

